% Lösungen Punkte und Strecken im Koordinatensystem · Dreiecke nachweisen · Prüfungs-Fokus Abitur GK (zusammenbau v0.6)
\einheitenkopf{Einheit 3 · Dreiecke nachweisen}

\erg{1}{a) $\overrightarrow{DC}$}
\erg{2}{a) $|\overrightarrow{CA}| = |\overrightarrow{CB}| = \sqrt{21}$, $|\overrightarrow{AB}| = \sqrt{20}$: gleichschenklig, Schenkel $CA$ und $CB$, Basis $AB$ \quad b) $|\overrightarrow{CA}| = |\overrightarrow{CB}| = \sqrt{14}$, $|\overrightarrow{AB}| = \sqrt{20}$: gleichschenklig, Schenkel $CA$ und $CB$, Basis $AB$ \quad c) $|\overrightarrow{RP}| = |\overrightarrow{RQ}| = \sqrt{29}$, $|\overrightarrow{PQ}| = \sqrt{20}$: gleichschenklig, Schenkel $RP$ und $RQ$, Basis $PQ$ \quad d) $|\overrightarrow{CA}| = |\overrightarrow{CB}| = \sqrt{14}$, $|\overrightarrow{AB}| = \sqrt{24}$: gleichschenklig, Schenkel $CA$ und $CB$, Basis $AB$ \quad e) $|\overrightarrow{CA}| = |\overrightarrow{CB}| = \sqrt{29}$, $|\overrightarrow{AB}| = 6$: gleichschenklig, Schenkel $CA$ und $CB$, Basis $AB$ \quad f) $|\overrightarrow{AB}| = |\overrightarrow{BC}| = |\overrightarrow{CA}| = \sqrt{8}$: gleichseitig. \quad g) $|\overrightarrow{AB}| = \sqrt{18}$, $|\overrightarrow{BC}| = |\overrightarrow{CA}| = 5$: nicht gleichseitig, nur gleichschenklig mit der Basis $AB$. \quad h) $|\overrightarrow{AB}| = |\overrightarrow{BC}| = |\overrightarrow{CA}| = \sqrt{32}$: gleichseitig. \quad i) $|\overrightarrow{CA}| = |(-2 | -2 | -p)| = \sqrt{8 + p^2}$ und $|\overrightarrow{CB}| = |(2 | 2 | -p)| = \sqrt{8 + p^2}$: gleich für jedes $p$, Basis $AB$. \quad j) $|\overrightarrow{CA}| = |\overrightarrow{CB}| = \sqrt{(3 - p)^2 + 10}$, denn $\overrightarrow{CA} = (3 - p | -3 | 1)$ und $\overrightarrow{CB} = (3 - p | 3 | -1)$ unterscheiden sich nur im Vorzeichen zweier Koordinaten. \quad k) $|\overrightarrow{CA}| = |\overrightarrow{CB}| = \sqrt{8 + (2 - p)^2}$ für jedes $p$; Basis $AB$. \quad l) Mit dem Ursprung $O$ ist $OABC$ ein Quadrat (Seiten auf den Achsen bzw. parallel dazu, alle gleich lang); $ABC$ ist seine Hälfte mit dem rechten Winkel bei $B$ (nicht bei $A$ oder $C$) und $|AB| = |BC|$. Flächeninhalt $\frac12 \cdot 9 \cdot 9 = 40{,}5$.}
\erg{3}{a) $|\overrightarrow{AB}| = |\overrightarrow{AC}| = \sqrt{34}$, Basis $BC$ (nicht Schenkel); $E$ und $F$ haben dieselbe dritte Koordinate $2$, die Grundfläche liegt in der $x_1x_2$-Ebene: $EF$ ist parallel zu ihr. \quad b) $|\overrightarrow{AB}| = |\overrightarrow{AC}| = 5$, Basis $BC$ (nicht Schenkel); $G$ und $H$ haben dieselbe dritte Koordinate $5$, die Grundfläche liegt in der Ebene $x_3 = 3$: $GH$ ist parallel zu ihr. \quad c) $|\overrightarrow{AB}| = |\overrightarrow{AC}| = \sqrt{29}$, Basis $BC$ (nicht Schenkel); $P$ und $Q$ haben dieselbe dritte Koordinate $3$, die Grundfläche liegt in der $x_1x_2$-Ebene: $PQ$ ist parallel zu ihr. \quad d) $M$ ist der Mittelpunkt von $DE$; $\overrightarrow{FD} \circ \overrightarrow{FE} = 0$, also hat das Dreieck bei $F$ einen rechten Winkel. Nach Thales liegt $F$ auf dem Kreis um $M$ durch $D$ und $E$. \quad e) $\overrightarrow{FD} \circ \overrightarrow{FE} = 0$: rechter Winkel bei $F$. Nach Thales liegt $F$ auf dem Kreis mit dem Durchmesser $DE$, sein Mittelpunkt ist die Mitte von $DE$ – drei Abstände muss man nicht rechnen. \quad f) $\overrightarrow{RP} \circ \overrightarrow{RQ} = 0$: rechter Winkel bei $R$. Nach Thales liegt $R$ auf dem Kreis mit dem Durchmesser $PQ$, sein Mittelpunkt ist die Mitte von $PQ$ – drei Abstände muss man nicht rechnen. \quad g) Symmetrie zu $AB$: $M(3 | m)$; $|MA| = |MC|$: $3^2 + m^2 = (9 - m)^2$ ergibt $m = 4$: $M(3 | 4)$ (nicht der Schwerpunkt $(3 | 3)$) \quad h) Symmetrie zu $AB$: $M(2 | m)$; $|MA| = |MC|$: $4^2 + m^2 = (8 - m)^2$ ergibt $m = 3$: $M(2 | 3)$ \quad i) Symmetrie zu $AB$: $M(4 | m)$; $|MA| = |MC|$: $3^2 + (m - 1)^2 = (10 - m)^2$ ergibt $m = 5$: $M(4 | 5)$}
\erg{4}{a) $C$ am Mittelpunkt $M(3 | 3 | 3)$ von $AB$ spiegeln: $D = A + B - C = (1 | 3 | 5)$; $ADBC$ ist dann eine Raute. \quad b) $C$ am Mittelpunkt $M(2 | 2 | 1)$ von $AB$ spiegeln: $D = A + B - C = (0 | 1 | -1)$; $ADBC$ ist dann eine Raute. \quad c) $|\overrightarrow{MA}| = 3$; $B = 2M - A = (-1 | 5 | 0)$ (Gegenpunkt); $C = M + (1 | 2 | 2) = (2 | 5 | 3)$, denn $(1 | 2 | 2)$ steht senkrecht auf $\overrightarrow{MA}$ und hat die Länge $3$ – dann ist $|CA| = |CB|$. \quad d) $|\overrightarrow{MA}| = 3$; $B = 2M - A = (6 | 4 | 3)$ (Gegenpunkt); $C = M + (1 | -2 | -2) = (5 | 0 | 2)$, denn $(1 | -2 | -2)$ steht senkrecht auf $\overrightarrow{MA}$ und hat die Länge $3$ – dann ist $|CA| = |CB|$.}
